\chapter{Proposed Approach}
\section{Notation and Experimental Isolation}
Let $D=\{(x_i,y_i)\}_{i=1}^{N}$ be a dataset with $C$ classes, represented by
numeric feature vectors and reference labels. Stratified splitting produces
training, validation, and test partitions in approximate proportions
60\%, 20\%, and 20\%. The split seed is shared across all conditions belonging
to a dataset--seed pair. The validation partition is reserved but not used
for tuning or training in this study. Test labels are never corrupted or used
to form cleaning scores.

Within the training partition, corruption produces observed labels
$\widetilde y_i$ and a Boolean mask $z_i=\mathbf{1}[\widetilde y_i\ne y_i]$.
The detectors observe $(x_i,\widetilde y_i)$, not $y_i$ or $z_i$.
Reference training labels are used only for the clean baseline, corruption
construction, and detection evaluation. This separation prevents a detector
from exploiting the answer it is supposed to recover.

\section{Controlled Label Corruption}
\subsection{Symmetric Noise}
For a nominal corruption rate $\eta$ and $n$ training rows, sample
$k=\operatorname{round}(\eta n)$ distinct indices without replacement. Replace
each selected label with a uniformly chosen member of the other $C-1$ classes.
Every selected label changes. In expectation, its transition probabilities
are described by
\begin{equation}
 T_{ab}=\Pr(\widetilde y=b\mid y=a)=
 \begin{cases}1-\eta,&a=b,\\ \eta/(C-1),&a\ne b.\end{cases}
\end{equation}
The implementation fixes the total number of changes rather than sampling
independent Bernoulli flips. Finite-sample transition frequencies therefore
need not equal this idealized matrix exactly.

\subsection{Cyclic Class-Dependent Noise}
For sorted class labels $c_0,\ldots,c_{C-1}$, define
$g(c_j)=c_{(j+1)\bmod C}$. Select
$k_j=\operatorname{round}(\eta n_j)$ rows from each original class and replace
their labels using $g$. Class membership is computed before any labels are
changed, so an example cannot be flipped twice as the class loop proceeds.
The nominal transition matrix has diagonal mass $1-\eta$ and off-diagonal
mass $\eta$ at $g(c_j)$.

This model is class-dependent in the destination of a flip, not in a
deliberately different error rate for each class. It is asymmetric for more
than two classes, but binary cyclic and symmetric corruption both send a
selected label to the only alternative class. Their binary implementations
still differ in global versus per-class sample selection and rounding.
The term ``cyclic'' is used to avoid implying a realistic annotator-confusion
process. With per-class rounding, $\sum_j k_j$ can differ slightly from
$\operatorname{round}(\eta n)$; exact counts are saved in the result matrix.

\section{Training-Only Out-of-Fold Probabilities}
For each noisy dataset and base classifier, use shuffled stratified five-fold
cross-validation within the training partition. Each training row receives
a probability vector from a model fitted without that row:
\begin{equation}
 p_i(c)=f^{(-v(i))}(x_i)_c,
\end{equation}
where $v(i)$ denotes its fold. The seed for fold shuffling is the experiment
seed. If a class has fewer than five rows, the implementation reduces the
fold count to the smallest class count; fewer than two rows raises an error.
The canonical data support the requested five folds.

Logistic regression's scaler is part of the fitted pipeline. Consequently,
each fold estimates its own scaling parameters using only its fitting rows.
The tree models use the numeric features without standardization. All four
detectors share the same out-of-fold probability matrix, avoiding additional
fitting variation or repeated cross-validation work. The same base model
family supplies both the detector and the downstream classifier; this is not
a cross-model transfer study.

\section{Suspiciousness Scores}
Define $\widehat y_i=\arg\max_c p_i(c)$, with the first sorted class chosen in
an exact probability tie. Let $p_i^{(1)}$ and $p_i^{(2)}$ be the largest and
second-largest predicted probabilities. Higher values of all four scores
indicate greater suspicion:
\begin{align}
 s_i^{\mathrm{conf}} &= 1-p_i(\widetilde y_i),\\
 s_i^{\mathrm{dis}} &= \mathbf{1}[\widehat y_i\ne\widetilde y_i],\\
 s_i^{\mathrm{margin}} &= 1-(p_i^{(1)}-p_i^{(2)}),\\
 s_i^{\mathrm{comb}} &= 0.5s_i^{\mathrm{conf}}+0.3s_i^{\mathrm{dis}}+0.2s_i^{\mathrm{margin}}.
 \label{eq:combined}
\end{align}
All component scores lie in $[0,1]$. The nonnegative weights sum to one and
were fixed in the project design; they were not optimized using the test
results. Confidence conditions directly on the observed label. Disagreement
discards probability magnitude and creates many ties. Margin measures general
uncertainty without asking whether the observed label is the uncertain class.
That difference is central to interpreting the results.

\subsection{Why the Binary Combination Adds Little Ranking Information}
For two classes, let $s=1-p_i(\widetilde y_i)$. Away from a probability tie,
substitution into Eq.~\eqref{eq:combined} gives
\begin{equation}
 s^{\mathrm{comb}}=
 \begin{cases}
  0.9s,&s<0.5,\\
  0.1s+0.7,&s>0.5.
 \end{cases}
 \label{eq:binary}
\end{equation}
Both branches increase with $s$, and the second branch is above the first.
Thus the combined score preserves the confidence ordering except for the
handling of exact ties. At $s=0.5$, the disagreement term can depend on which
class wins the argmax tie. Matching binary performance is therefore expected
and is not strong evidence that the combination contributes a new signal.
For $C>2$, the top-two margin is not determined by the observed-label
probability, so the ranking can change. The multiclass results are the more
informative component comparison.

\section{Selection and Cleaning}
For every detector, select the $k=\operatorname{round}(\eta n)$ highest scores.
Python's nearest-even rounding is used. Sorting is stable, ascending, and the
last $k$ indices are selected. Exact ties at the cutoff therefore favor later
rows in the current training-row order. That deterministic rule is especially
consequential for the binary-valued disagreement score.

The use of $\eta$ is an \emph{oracle selection budget}: the fraction of labels
injected by the experiment is assumed known. The mask $z_i$ is not used for
selection, but this budget assumption is more favorable than real deployment
with an unknown noise level. For symmetric noise, the number of selected and
corrupted rows is identical by construction. Detection precision and recall
then coincide. Under cyclic corruption they may differ slightly because the
number of injected errors is rounded separately within each class.

\textbf{Removal} discards selected examples and fits a fresh classifier on
the remaining noisy-labeled training set. \textbf{Relabeling} keeps every
feature row and replaces selected labels by $\widehat y_i$, then fits a fresh
classifier. A selected row already predicted as its observed label remains
unchanged after relabeling. Relabeling uses one model family's out-of-fold
predictions, not an ensemble consensus or a second round of label recovery.

Two references complete the comparison: a fresh model trained on original
training labels, and a fresh model trained on all noisy training labels.
The original code identifier \texttt{clean\_upper\_bound} denotes the former;
it is a reference condition, not a mathematical upper bound on test accuracy.

\section{Experiment Algorithm and Cost}
For each dataset, noise type, rate, model, and seed, the runner:
\begin{enumerate}
\item constructs the fixed training/validation/test split and training corruption;
\item fits and evaluates clean and noisy baseline models;
\item computes one shared out-of-fold probability matrix on noisy training data;
\item evaluates each detector against the injected-corruption mask;
\item removes or relabels its selected rows, retrains, and evaluates on the same test set;
\item writes the ten result rows and an atomic training-provenance checkpoint.
\end{enumerate}
With five folds, each combination requires two baseline fits, five
out-of-fold fits, and eight cleaned-model fits: 15 top-level estimator fits.
Across 384 combinations this is 5,760 fits, excluding subsequent diagnostic
retraining and reproducibility checks. A forest fit itself builds 200 trees.
``Lightweight'' describes simple scores on conventional CPU estimators; it
does not mean the full benchmark is cost-free or that runtime superiority was
measured. The study does not provide a controlled timing comparison.
